\chapter{Certificats de précision arbitraire et échantillons décimaux finis reproductibles}
\label{app:desarrollos-10000-rev9}

Cet appendice imprime les sorties complètes de l'exécution décrite dans
\cref{sec:truncamiento-10000-pi-e-phi-rev9,sec:alpha-t1-10000-rev9}. Les
sauts de ligne sont exclusivement typographiques. Chaque développement constitue
un témoin fini de la loi coinductive de précision arbitraire démontrée dans
la Partie III ; il ne définit pas la constante ni ne limite la profondeur du générateur.

\begingroup
\lstset{basicstyle=\ttfamily\tiny,breaklines=true,breakatwhitespace=false,
  frame=single,columns=fullflexible,keepspaces=true,numbers=none}

\section{Développement coinductif de \(\pi\) et certificat de troncature compatible}
\lstinputlisting{generated/decimals_10000_rev9/pi_10000_decimales_96col.txt}

\section{Développement coinductif de \(e\) et certificat de troncature compatible}
\lstinputlisting{generated/decimals_10000_rev9/e_10000_decimales_96col.txt}

\section{Développement coinductif de \(\varphi\) et certificat de troncature compatible}
\lstinputlisting{generated/decimals_10000_rev9/phi_10000_decimales_96col.txt}

\section{Développement de \texorpdfstring{\(\alpha_{\rm HMT}\)}{alpha HMT}}
\lstinputlisting{generated/decimals_10000_rev9/alpha_T1_10000_decimales_96col.txt}
\endgroup
